\section{Discussion and Limitations}\label{sec:discussion}

\paragraph{Which rival the speedup beats} Proposition~\ref{prop:advantage} compares the quantum interface with classical algorithms that learn about the network only by evaluating the flow of candidate plans. Exhaustive enumeration of $\Fk$ is the best algorithm of this kind up to a factor of two, and the simulations confirm the separation that the proposition predicts (Section~\ref{sec:results-scaling}). The comparison says nothing about algorithms that exploit the structure of the network, and the MILP is one of them. On the instances of E1 it was several orders of magnitude faster than both searches (Section~\ref{sec:results-classical}), so the query-level speedup holds only relative to exhaustive search.

\paragraph{When the speedup matters} When $k$ grows with the network, the quantum algorithm still needs $\Theta(\sqrt{N})$ oracle calls, which is exponential in $k$. Proposition~\ref{prop:advantage} is thus a quadratic speedup over exhaustive search and not an efficient algorithm for an NP-hard problem.

\paragraph{Cost of the oracle} The proposition counts oracle calls of polynomial size, but the polynomial matters. The arithmetic realization of Section~\ref{sec:interface-oracle}, the only one we built for general instances, already needs hundreds to thousands of two-qubit gates on networks with at most 14 arcs (Table~\ref{tab:resources}). It also searches the larger space of pairs $(x,\alpha)$ and can beat enumeration of $\Fk$ in query count only when $2^{n-2}<N$, which held on 7 of the 9 instances of E2. A reversible maximum flow circuit, which the speedup over $\Fk$ requires, would be larger still, and Proposition~\ref{prop:advantage} currently relies on an existence argument for it rather than on an explicit construction. Babbush \emph{et al.}~\cite{babbush2021} estimate that quadratic speedups need very long runs to offset the overhead of error correction on early fault-tolerant hardware. For the MFNIP, where every query contains a maximum flow computation, the crossover with classical enumeration would appear only for very large $N$, and the crossover with the MILP may never occur.

\paragraph{General costs} The interface depends on a uniform superposition of the feasible plans. With unit costs, Dicke states provide this superposition. With arbitrary integer costs, it becomes a uniform superposition over the solutions of a knapsack constraint, for which we know no efficient preparation. The provable speedup therefore holds only for the cardinality variant.

\paragraph{What the hardware experiment shows} The truth-table oracle used on the processor marks a plan found beforehand by classical means, and its cost grows exponentially with $m$. The experiment shows that the interface, consisting of the preparation, the oracle, and the reflection, amplifies the optimal plan on a three-arc instance. It does not demonstrate a speedup and would not scale. At the optimal number of iterations, the transpiled interface circuit had fewer two-qubit gates than the hypercube circuit (29 against 39) and a higher frequency of the optimum on the noise model (0.898 against 0.833). On the processor, however, the two frequencies were practically equal (0.815 and 0.818, Table~\ref{tab:hardware}), so the higher ideal probability of the interface variant did not carry over to the hardware. The noise model predicted the hypercube circuit well but overestimated the interface circuit. Because we defined these comparisons after the run, a new hardware run with the comparison against the noise model registered in advance is needed to confirm them. Noise also produced plans outside the budget, which the classical loop must discard after evaluating them.

\paragraph{Threats to validity} The scaling experiment simulates the algorithm at the query level, so it verifies the query complexity of the interface and the correctness of the control loop but not the behavior of a physical oracle. The link between the simulator and the circuits was verified only on small instances (Section~\ref{sec:method-validation}). The instances come from two synthetic generators and are small. The layered family has a fixed number of arcs for each width, so $N$ grows together with $k$ and $m$. The number of optimal plans $M$ also varies across instances, which is why we also fit against $N/M$. Finally, the literature search covered titles and abstracts in four sources and may have missed work indexed elsewhere.
